% Shared preamble for the printable mini lessons (1.4 onward).
%   \documentclass[11pt]{article}
%   \input{../preamble}
% Build the key with:  pdflatex -jobname=<name>_KEY "\def\iskey{}\input{<name>}"
% Figures shared with the web page live in figures/Unit_X/<name>.tex and are
% pulled in with \fig{Unit_X/<name>}; the standalone package drops their own
% preamble, so their macros come from figures/preamble.tex, loaded here.
\input{../../figures/preamble}
\usepackage[margin=0.75in]{geometry}
\usepackage{array}
\usepackage{enumitem}
\usepackage{amsmath}
\usepackage{cancel}
\usepackage[helvet]{sfmath}  % math in the same sans font as the text
\usepackage{standalone}
\renewcommand{\ttdefault}{pcr}  % Courier: a real Type 1 typewriter font
\pagestyle{empty}
\setlength{\parindent}{0pt}
\setlist[enumerate,1]{leftmargin=1.5em,itemsep=1em}
\setlist[enumerate,2]{label=(\alph*),leftmargin=2em,itemsep=2.2em}

\newif\ifkey
\ifdefined\iskey \keytrue \fi
\colorlet{keyred}{red!75!black}
\newcommand{\ans}[1]{\ifkey\leavevmode{\color{keyred}\let\frac\dfrac #1}\fi}  % full-size fractions in answers
% vertical work space on the student copy, less on the key
\newcommand{\work}[1]{\ifkey\vspace{0.6em}\else\vspace{#1}\fi}
\newcommand{\fig}[1]{\input{../../figures/#1}}
% Title line: \lessontitle{Quartiles, IQR, and Box Plots}
% A long title gets the Name line on its own line above it.
\newlength{\titlewidth}
\newcommand{\lessontitle}[1]{%
  \settowidth{\titlewidth}{\Large\bfseries #1\ifkey\ KEY\fi}%
  \ifdim\titlewidth>4.3in
    \hfill Name \underline{\hspace{2in}}\par\vspace{0.5em}
    {\Large\bfseries #1\ifkey\ \textcolor{keyred}{KEY}\fi}\par\vspace{0.6em}
  \else
    {\Large\bfseries #1\ifkey\ \textcolor{keyred}{KEY}\fi}\hfill Name \underline{\hspace{2in}}\par\vspace{0.6em}
  \fi}
\newcommand{\partner}{\par\vspace{0.4em}Work with a partner. Compare answers when you finish each problem.\par}

% ---------- shared pieces added for AI and You ----------
% no hyphenated words anywhere: students read in a second language
\hyphenpenalty=10000 \exhyphenpenalty=10000
\usepackage{pifont}   % \ding{34} scissors
\newcommand{\blank}[1]{\underline{\hspace{#1}}}
% ragged table columns: no hyphenated words for students reading in a second language
\newcolumntype{R}[1]{>{\raggedright\arraybackslash}p{#1}}
% An answer written on a blank line: \ansline{width}{answer}
\newcommand{\ansline}[2]{\ifkey\underline{\makebox[#1][l]{\,\textcolor{keyred}{#2}}}\else\blank{#1}\fi}
% Unit 1 class timeline: a card to cut out and post on the wall.
%   \timelinecard{1950}{Alan Turing asks ...}
\newcommand{\timelinecard}[3][8.3]{%
  \begin{tikzpicture}
    \pgfmathsetmacro{\tw}{#1-0.8}
    \draw[dashed, line width=0.8pt, rounded corners=3pt] (0,0) rectangle (#1,4.7);
    \node[anchor=north west, font=\huge\bfseries] at (0.3,4.45) {#2};
    \node[anchor=north west, text width=\tw cm, align=left, font=\normalsize] at (0.3,3.45) {#3};
    \node[anchor=south east, font=\footnotesize, text=black!55] at (#1-0.2,0.15) {Draw a picture here.};
  \end{tikzpicture}}
% three cards across: \timelinecard[5.6]{...}{...}
\newcommand{\cuthere}{\noindent{\Large\ding{34}}\,{\small\itshape Cut out the cards. Add them to the class timeline.}\par\vspace{0.3em}}
